\section{图的语言与度数}
\label{sec:04-Discrete-Mathematics/03-Graph-Theory/01}

你手里有一份社交网络数据：500 个人，有些人是好友，有些人不是。你想回答几个最朴素的问题——谁的朋友最多？有没有人完全孤立？整个网络是连成一片还是碎成几个小圈子？你不可能对着 500 行的名单逐对查关系。

你需要一种语言，把``谁认识谁''这种关系翻译成能计算的数学对象。画 500 个点和连线会糊成一团，但数学语言不在乎画面是否美观——它只要求你明确声明：哪些对象存在，哪些对之间有连接。

\subsection{第一件事：宣布边代表什么}

图 \(G=(V,E)\) 由顶点集 \(V\) 和边集 \(E\) 组成。顶点是你关心的对象——人、城市、任务、网页。边是你声明的关系——好友、道路、依赖、超链接。

但``关系''有不同种类。社交平台上的好友通常是双向的：A 是 B 的好友意味着 B 也是 A 的好友。这种对称关系用无向边表示，边是顶点的无序对 \(\{u,v\}\)。网页之间的超链接是单向的：A 页面链接到 B 页面，B 未必链接回 A。这时你只能用有向边，边是有序对 \((u,v)\)，从尾部指向头部。道路的长度或通信的延迟则需要在边上标注数值，称为权重。

这些不是细枝末节。你把有向关系当成无向关系建模，会得出``互相可达''的错误结论；你忽略了权重，就分不出``两步短路''和``绕城一圈''的差别。

还有两条约定需要在开始时说明。简单图不允许自环（顶点连向自己的边）和平行边（同一对顶点间多条边）。但如果你在建模一个路网，两个路口之间可能确有两条平行的单行道，这时多重图才是诚实的模型。一个结论是否依赖``简单图''这个前提，应该在陈述时标明。

最后，图的画法不是图。两个看起来完全不同的布局——一个散乱、一个对称——只要顶点之间的邻接关系完全一致，它们就代表同一个图（同构）。画图时边的交叉也不代表有交点，除非你在交叉处画了一个顶点。

\subsection{度数：给每个顶点的连接数单子}

无向图中，顶点 \(v\) 的度 \(\deg(v)\) 就是与 \(v\) 相连的边数。这个数足够简单：拉一张好友列表，数一下有几个人。

现在把所有人的度数加起来。每条边有两个端点，在求和时被各数一次，所以：

\[
\sum_{v\in V}\deg(v)=2|E|.
\]

这条关系被称为握手定理——聚会上所有人握手次数的总和等于握手次数的两倍，因为每次握手涉及两个人。

从这个等式能立刻推出一个不显然的事实：图中奇度顶点的个数必为偶数。理由是总和 \(2|E|\) 是偶数，而若干个整数之和为偶数时，其中的奇数必须成对出现（奇数加奇数是偶数，但单独的奇数会让总和变奇）。这句话只用了度数和边的全局等式，没翻开任何一个顶点的邻居名单，就已经对全图的度数奇偶分布做出了限制。

有向图的度数拆成入度和出度：入度 \(\deg^-(v)\) 是从别的顶点指向 \(v\) 的边数，出度 \(\deg^+(v)\) 是从 \(v\) 指向别处的边数。每条有向边为一个顶点贡献一次出度、为另一个贡献一次入度，因此：

\[
\sum_{v}\deg^+(v)=\sum_{v}\deg^-(v)=|E|.
\]

\subsection{子图：从大图里切出局部}

有时你只关心一部分顶点及它们之间的连接。子图可以从原图删除顶点、删除边，或两者都删。但有一个更精确的操作：给定顶点子集 \(S\subseteq V\)，由 \(S\) 诱导的子图保留原图中两端都在 \(S\) 内的全部边——不增不减。当你想要断言``这群对象内部具有某种性质''时，诱导子图是你真正在讨论的对象；随便删几条边再讨论就会漏掉重要信息。

简单图的补图 \(\overline{G}\) 在同一个顶点集上，把原来没有的边全部补上、原来有的边全部删去。\(G\) 和 \(\overline{G}\) 的边恰好瓜分了完全图 \(K_n\) 的全部 \(\binom{n}{2}\) 条边：

\[
|E(G)|+|E(\overline{G})|=\binom{n}{2}.
\]

这个等式在你不方便直接数 \(G\) 的边时提供了一个替代路径——去数补图的边。

\subsection{选什么表示，取决于你要做什么操作}

同样的抽象图，计算机里可以用完全不同的数据结构存储。

邻接表为每个顶点维护一个邻居列表。遍历一个顶点的所有邻居只需扫描这个列表，时间与度数成正比。适合稀疏图（边数远小于 \(\binom{n}{2}\)），也是 BFS 和 DFS 的自然搭档。

邻接矩阵是一个 \(n\times n\) 的方阵 \(A\)，\(A_{ij}=1\) 当且仅当顶点 \(i\) 和 \(j\) 之间有边（无向图的矩阵是对称的）。它占 \(O(n^2)\) 空间，不论边有多少，但换来了一个强大的代数性质：\((A^k)_{ij}\) 恰好等于从 \(i\) 到 \(j\) 的长度为 \(k\) 的游走数。每左乘一次 \(A\)，相当于在当前游走的末端延长一步——这可以用归纳法直接验证。如果你需要回答``任意两点间长度不超过 \(k\) 的路径是否存在''，矩阵乘法是你的工具。

边表只是把所有边列成一个清单。它适合按边权排序的算法——比如 Kruskal 最小生成树算法每次取出最轻的边，边表排序一次即可。

表示法改变的是算法的时间复杂度和代码实现，不改变图的数学性质。选哪一种，取决于你最频繁的操作是``枚举邻居''还是``查询任意两点是否相邻''还是``按边权排序''。

下一步看路径和连通性：一旦你有了邻接关系这张网，如何追踪影响力从一点传播到另一点。
